\chapter[1943 --- Doisy]{1943: Henrik Carl Peter Dam and Edward Adelbert Doisy}
\label{ch:1943_Dam}

\section*{Main Contribution}

\textit{Discovered vitamin K and its essential role in blood coagulation, explaining why some newborns bleed excessively and how to prevent it.}\cite{nobel-1943,lecture-1943}

\section{Official Citation}

for his discovery of the role of vitamin K in the physiology of coagulation of the blood

\section{Background and Historical Context}

\subsection{Henrik Carl Peter Dam}

Henrik Carl Peter Dam was born on February 21, 1895, in Copenhagen, Denmark. He grew up in a cultured and educated family; his father, Alfred Dam, was a successful businessman, and his mother, Emilie Peterson, came from a family of teachers. Dam received his early education in Copenhagen, where he demonstrated a strong aptitude for the natural sciences. He enrolled at the Copenhagen Polytechnic Institute (now the Technical University of Denmark) in 1918, studying chemistry and agricultural science, and graduated in 1923 with a degree in chemical engineering.

After graduation, Dam joined the Biochemical Institute of the Royal Veterinary and Agricultural College in Copenhagen, where he began research on the role of vitamins in animal nutrition. His early work focused on vitamin A, and he made important contributions to understanding its distribution and function. However, it was his work on a previously unrecognized vitamin---later designated vitamin K---that would bring him international recognition and the Nobel Prize.

The scientific context of the 1930s was one of rapid progress in vitamin research. The discovery and characterization of vitamins A, B, C, and D had already earned Nobel Prizes for several researchers, and there was intense interest in identifying additional dietary factors essential for health. Dam's work on vitamin K emerged from his studies on cholesterol metabolism in chickens, which led to the observation that a specific dietary factor was necessary for normal blood clotting.

\subsection{Edward Adelbert Doisy}

Edward Adelbert Doisy was born on November 13, 1893, in Hume, Illinois, USA. He was raised by his mother, who died when he was just seven years old, and was subsequently cared for by his aunt. Despite these early difficulties, Doisy excelled academically and attended the University of Illinois, where he earned his bachelor's degree in 1914 and his master's degree in 1916. He then pursued graduate studies at Harvard University, where he received his Ph.D. in biochemistry in 1920 under the supervision of Otto Folin.

Doisy joined the faculty of Saint Louis University School of Medicine in 1919, where he would spend the next four decades conducting research on the sex hormones and vitamins. His early work focused on the isolation and characterization of estrin (now called estrogen), the female sex hormone, and he was the first to isolate this compound in pure form. This work established his reputation as one of the leading biochemists in the United States and laid the groundwork for his subsequent work on vitamin K.

\section{The Discovery/Contribution}

The Nobel Prize was awarded to Dam and Doisy jointly: to Dam ``for his discovery of the role of vitamin K in the physiology of coagulation'' and to Doisy ``for his discovery of the chemical nature of vitamin K.'' Together, their work identified, characterized, and elucidated the function of vitamin K---a discovery that has had significant implications for clinical medicine.

\subsection{Henrik Dam's Discovery of Vitamin K}

Dam's discovery of vitamin K emerged from his studies on cholesterol metabolism in chickens. In 1929, while conducting experiments in which he fed chicks a diet supplemented with cholesterol from which the lipid-soluble vitamins had been extracted, Dam observed that the animals developed a bleeding disorder characterized by prolonged clotting time. Despite the fact that their diet contained adequate amounts of vitamins A and D, and that they showed no signs of scurvy (vitamin C deficiency), the chicks developed spontaneous hemorrhages and died if the diet was continued.

Dam initially hypothesized that the cholesterol preparation contained a toxic factor that was causing the clotting abnormality. However, systematic investigation revealed that the problem was not the presence of a toxin but rather the absence of an essential dietary factor. When Dam added certain foods---particularly green leafy vegetables and certain cereals---to the deficient diet, the clotting abnormality was corrected. Through painstaking fractionation of these foods, Dam was able to isolate the active factor, which he designated ``vitamin K'' (from the German word ``Koagulation'').

Dam's work on the physiological role of vitamin K was equally important. He demonstrated that vitamin K is essential for the synthesis of prothrombin (now known as coagulation factor II), a key protein in the blood clotting cascade. In the absence of vitamin K, the liver produces an inactive form of prothrombin, resulting in impaired clotting and a tendency to bleed excessively. This finding established the fundamental link between vitamin K and blood coagulation and provided the basis for understanding the vitamin's clinical significance.

\subsection{Edward Doisy's Isolation and Characterization of Vitamin K}

While Dam was elucidating the biological role of vitamin K, Doisy was independently working on its chemical characterization. Using a combination of chemical extraction techniques and bioassays, Doisy and his collaborators isolated vitamin K from various sources, including alfalfa meal and fish meal. In 1939, Doisy published the chemical structure of vitamin K1 (phylloquinone), which was shown to be a substituted 2-methyl-1,4-naphthoquinone with a long isoprenoid side chain. Shortly thereafter, Doisy also identified vitamin K2 (menaquinone), a related compound with a longer side chain that is produced by bacteria in the intestine.

Doisy's determination of the chemical structure of vitamin K was a tour de force of analytical chemistry. The compound was present in extremely small quantities in natural materials, and isolating sufficient quantities for structural analysis required processing thousands of kilograms of raw material. Doisy's success in this challenging endeavor was a testament to his chemical expertise and his determination to solve the problem.

The structural characterization of vitamin K opened the door to its total synthesis, which Doisy and others achieved within a few years of the initial isolation. The ability to synthesize vitamin K made possible the production of the vitamin in pharmaceutical form, which has been used ever since in the prevention and treatment of bleeding disorders.

\section{Impact on Modern Medicine}

The discovery of vitamin K has had far-reaching consequences for clinical medicine, affecting fields as diverse as pediatrics, surgery, hematology, and preventive medicine.

\subsection{Neonatal Vitamin K Deficiency}

One of a major clinical applications of Dam's and Doisy's work is the prevention of vitamin K deficiency in newborn infants. Newborns are particularly vulnerable to vitamin K deficiency because the vitamin crosses the placenta poorly and is present in low concentrations in breast milk. Without prophylaxis, approximately 1 in 300 infants develops vitamin K deficiency bleeding (VKDB), formerly known as hemorrhagic disease of the newborn, which can cause life-threatening intracranial hemorrhage.

The routine prophylactic administration of vitamin K to all newborns---a practice that began in the 1960s and is now standard of care in most developed countries---has virtually eliminated VKDB. This simple intervention, made possible by the discoveries of Dam and Doisy, saves thousands of lives each year and represents one of the great successes of preventive medicine in the twentieth century.

\subsection{Surgical and Medical Applications}

Vitamin K is essential for patients undergoing surgery, as it ensures normal blood clotting and reduces the risk of perioperative bleeding. Patients who are at risk of vitamin K deficiency---including those with liver disease, biliary obstruction, or malabsorption syndromes---are routinely given vitamin K supplements before and after surgical procedures to optimize their coagulation status.

In emergency medicine, vitamin K (phytonadione) is used to reverse the effects of warfarin, an anticoagulant drug that works by inhibiting the vitamin K-dependent clotting factors. When patients on warfarin therapy need emergency surgery or develop life-threatening bleeding, intravenous vitamin K can rapidly restore normal clotting function. This application represents a direct translation of the basic science discoveries of Dam and Doisy into clinical practice.

\subsection{Osteoporosis and Cardiovascular Disease}

Research over the past several decades has revealed that vitamin K plays important roles beyond blood coagulation. Vitamin K is essential for the carboxylation of osteocalcin, a protein involved in bone mineralization, and of matrix Gla protein, which helps prevent vascular calcification. These findings have led to interest in the potential role of vitamin K supplementation in the prevention and treatment of osteoporosis and cardiovascular disease, areas that are the subject of ongoing clinical research.

Epidemiological studies have suggested that higher dietary intake of vitamin K is associated with reduced risk of both osteoporotic fractures and cardiovascular disease, although the results of clinical trials have been mixed. Nevertheless, these emerging roles for vitamin K represent exciting new frontiers in nutrition and preventive medicine that build directly on the foundational work of Dam and Doisy.

\subsection{Anticoagulant Therapy}

The understanding of vitamin K's role in coagulation led directly to the development of warfarin and related anticoagulant drugs, which are among the most widely prescribed medications in the world. Warfarin works by inhibiting the enzyme vitamin K epoxide reductase, which is necessary for recycling the active form of vitamin K. By blocking this enzyme, warfarin reduces the production of functional clotting factors, thereby preventing the formation of dangerous blood clots in patients with conditions such as atrial fibrillation, deep vein thrombosis, and pulmonary embolism. The development of warfarin is a direct descendant of Dam and Doisy's research, and it remains a cornerstone of anticoagulant therapy more than seventy years after its introduction.

\section{Legacy}

The legacies of Henrik Dam and Edward Doisy are inextricably linked through their joint contribution to the discovery and characterization of vitamin K. Dam's identification of the vitamin and its role in coagulation, combined with Doisy's determination of its chemical structure, provided a complete picture of this essential nutrient that continues to inform clinical practice today.

Dam continued his research at the University of Copenhagen after receiving the Nobel Prize, studying the role of vitamins in lipid metabolism and making additional contributions to the understanding of vitamin E and other fat-soluble vitamins. He received numerous honors, including membership in the Royal Danish Academy of Sciences and Letters and honorary degrees from universities around the world. He died on September 17, 1976, in Copenhagen, at the age of 81.

Doisy remained at Saint Louis University for the rest of his career, continuing his research on hormones and vitamins and serving as chairman of the Department of Biochemistry for many years. He received the National Medal of Science in 1971 and numerous other honors. He died on October 23, 1986, at the age of 93.

Together, Dam and Doisy exemplify the power of collaborative scientific discovery. Working independently on different aspects of the same problem---Dam in Denmark, investigating the biological role of the vitamin, and Doisy in the United States, determining its chemical structure---they combined their efforts to produce a comprehensive understanding of vitamin K that has benefited millions of people around the world. Their story illustrates the international nature of scientific research and the importance of basic science in driving advances in clinical medicine.


\section{Modern Developments}

Dam and Doisy's vitamin K work has evolved into modern hematology. Vitamin K antagonists (warfarin) remain widely used anticoagulants. Direct oral anticoagulants (DOACs: apixaban, rivaroxaban) offer alternatives with fewer monitoring requirements. Understanding of vitamin K-dependent clotting factors has enabled development of recombinant clotting factors for hemophilia. Newer anticoagulants targeting factor XIa are in development, potentially offering anticoagulation without increased bleeding risk.


\section{Bibliography}

\begin{thebibliography}{9}
\bibitem{nobel-1943}
Nobel Prize Outreach, ``The Nobel Prize in Physiology or Medicine 1943,''\emph{NobelPrize.org}.\\
\url{https://www.nobelprize.org/prizes/medicine/1943/summary/}

\bibitem{lecture-1943}
Henrik Carl Peter Dam, ``Nobel Lecture,''\emph{NobelPrize.org}. Winner-authored primary source; downloadable from the official Nobel Prize archive.\\
\url{https://www.nobelprize.org/prizes/medicine/1943/dam/lecture/}
\end{thebibliography}
